\chapter{Introduction} \label{ch:introduction}

%% DRAFT (Claude, 2026-10-16): tense pass, Objectives, Research Questions lead-in,
%% Research Methods and Structure redrafted from the settled decisions in
%% dissertation/DECISIONS.md. To be read and rewritten in the author's own words.

This chapter sets out the context and motivation of this work, its objectives, hypothesis and research questions, the method followed, and how the rest of the document is organised.

\section{Context}

In recent years, Large Language Models (LLMs) have transformed natural language processing (NLP) \cite{openai2024gpt4technicalreport, yang2023harnessingpowerllmspractice, NEURIPS2022_b1efde53}, excelling in tasks like chatbots, summarization, and content analysis. However, they often struggle with domain-specific questions, producing inaccurate or irrelevant results, especially when data is outdated or sparse \cite{kandpal2023largelanguagemodelsstruggle, rag-survey}. This is particularly problematic in political discourse, where positions evolve over time.

Political party manifestos are crucial documents that outline a party's policy positions and ideological stances. These texts play a significant role in shaping public discourse and informing policy decisions \cite{manifestos-nz}. In the case of Portugal, like many other countries, understanding the nuances within these documents is crucial for informed civic participation \cite{marino-etal-2024-decoding}. However, the complexity and length of these documents can make them difficult for the average citizen and even researchers to navigate and fully understand.


\section{Motivation}

Citizens and researchers often struggle to navigate and compare political party programs, especially during elections. Existing tools either lack interpretability or require significant manual effort \cite{laver_benoit_garry, manifestos-nz}. This work addresses that gap with \sysname{}, a system that uses Retrieval-Augmented Generation (RAG) to retrieve and synthesise what the party manifestos say, in European Portuguese, a language in which research on this problem remains limited.

Moreover, political divisions are deepening, often fueled by a lack of accessible information \cite{eleitorado}. A tool that helps citizens better understand party positions could contribute to a more informed electorate and mitigate misinformation.

\section{Objectives}

The primary objective was to find out how well RAG retrieves and synthesises political content, using the manifestos of the 2024 Portuguese legislative election as the case. It broke down into three goals:

\begin{itemize}
    \item \textbf{RAG implementations:} survey RAG techniques and test which of them work for political manifestos.
    \item \textbf{Challenges of RAG:} identify what makes Portuguese manifestos hard for a RAG system.
    \item \textbf{Solutions to challenges:} apply and measure techniques that address those challenges, such as token-sized chunking, contextual retrieval, hybrid dense and lexical search, and reranking.
\end{itemize}

The work delivered four things:

\begin{itemize}
    \item \textbf{\sysname{}}, a working prototype that answers questions about the eight manifestos, citing the passages each claim comes from (Chapters~\ref{ch:proposition} and~\ref{ch:prototype-demonstration}).
    \item \textbf{A variant ladder}: every retrieval technique sits behind a feature flag, so each one can be switched on alone and its contribution measured (Section~\ref{subsec:variants}).
    \item \textbf{An evaluation harness and gold set}: 61 questions written and reviewed by hand, each anchored to the manifesto passages that answer it, scored with retrieval, routing and generation metrics (Sections~\ref{subsec:goldset} and~\ref{subsec:metrics}).
    \item \textbf{A decision history}: each component records the options considered, the choice made and, where it happened, why the choice was later reversed (Chapter~\ref{ch:proposition}).
\end{itemize}

A RAG-powered political quiz, to improve the quizzes that appear in every election campaign, was set as an optional goal. It was not built; it is discussed as future work (Chapter~\ref{ch:conclusion}).

\section{Research Hypothesis}

The main hypothesis of this project is:

“RAG-based retrieval and synthesis can effectively improve the accessibility and interpretability of Portuguese political manifestos.”

This hypothesis suggests that RAG implementations can accurately retrieve and summarize relevant political content, providing a structured overview of party positions.

\section{Research Questions} \label{sec:research-questions}

The hypothesis is too broad to test directly, so it is split into five questions that the evaluation harness can answer. Each one maps to a part of \sysname{}: RQ1 to indexing, RQ2 to retrieval, RQ3 to generation, and RQ4 and RQ5 to the system as a whole.

\begin{description}
    \item[RQ1 --- Representation.] How do representation choices --- measuring
    chunk size in the encoder's own tokens, and enriching each chunk with its
    section path and an LLM-written context --- affect retrieval quality over
    Portuguese electoral manifestos?

    \item[RQ2 --- Retrieval.] Which retrieval techniques (structured query
    analysis with routing, multi-query expansion with rank fusion, hybrid
    dense and lexical search, cross-encoder reranking) improve retrieval over
    a naive baseline, and how much does each contribute on its own?

    \item[RQ3 --- Generation.] Do retrieval improvements translate into
    better answers, measured as faithfulness to the retrieved manifesto
    content, relevance to the question, and correctness against a reference
    answer?

    \item[RQ4 --- Cost.] What do these improvements cost in latency, tokens
    and money, and which of them are worth that cost for a tool intended for
    public use?

    \item[RQ5 --- Retrieval versus scale.] How does \sysname{} compare with
    asking a language model the same question without retrieval, in answer
    quality and in resources consumed? Does a small model with the manifestos
    match or beat a stronger model without them?
\end{description}

\TODO{RQ2 names routing and multi-query expansion, whose rungs
(\texttt{routed}, \texttt{multiquery}) are not in the planned test run. Either
add them to that run (no code needed) or drop them from the wording. If the
web-search assistant rungs (\texttt{assistant\_*}) stay in the headline run,
add ``or a commercial assistant with web search'' to RQ5.}

Two questions were deliberately left out. Whether a hand-curated gold set can stand in for a human-authored benchmark is treated as methodology (Section~\ref{subsec:goldset}), not as a research question. And whether retrieval in European Portuguese reaches the quality reported for English cannot be measured without a parallel English corpus; it is discussed as a limitation (Chapter~\ref{ch:conclusion}).

\section{Research Methods}

This research followed an experimental and analytical approach to assess the effectiveness of RAG-based retrieval for Portuguese political manifestos. The plan had the following steps:

\begin{itemize}
    \item \textbf{Literature review:} a review of research on RAG pipelines, political text analysis, and the related challenges in NLP.
    \item \textbf{Dataset compilation:} collecting and preprocessing the 2024 Portuguese legislative manifestos.
    \item \textbf{Development:}
    \begin{itemize}
        \item \textbf{RAG experimentation:} testing RAG configurations, such as chunking strategies, hybrid search and reranking.
        \item \textbf{Prototype development:} building \sysname{}, a working prototype that retrieves and synthesises manifesto content.
    \end{itemize}
    \item \textbf{Testing and evaluation:}
    \begin{itemize}
        \item \textbf{Quantitative methods:} retrieval metrics (hit rate, recall, mean reciprocal rank, nDCG) against hand-picked gold passages, and claim-level answer metrics scored by an LLM judge.
        \item \textbf{Qualitative methods:} reading answers for errors the metrics cannot see, such as political framing passing from the sources into the answer (Section~\ref{subsec:grounding}).
    \end{itemize}
    \item \textbf{Refinement:} improving \sysname{} based on the evaluation results.
\end{itemize}

The plan held, but the process was not linear. The field moved fast while this work ran (Section~\ref{subsec:sota-recent}), and several early choices were revisited: the embedding model, the vector store, and the unit in which chunks are measured. Revisiting a choice when better evidence arrived was part of the method, not a departure from it. Chapter~\ref{ch:proposition} records each reversal and its reason, and the variant ladder keeps every technique measurable on its own, so a reversal is a measured change rather than a guess.

\TODO{AI-assistance disclosure, to confirm and rewrite: AI assistants were
used for code (first ChatGPT, later Claude Code) and for drafting text, which
the author reviewed and rewrote; the gold set's draft reference answers and
gold passages were proposed by Claude and reviewed by hand
(Section~\ref{subsec:goldset}). Check whether the university has a required
wording or a separate declaration for this.}

\section{Structure}

The Preamble, before this chapter, tells how the research question changed: from classifying political bias in Portuguese text to answering questions over party manifestos. The rest of the document is organised as follows:

\begin{itemize}
    \item \textbf{Chapter~\ref{ch:sota}, State of the Art,} surveys RAG, its recent developments, political text analysis and RAG evaluation.
    \item \textbf{Chapter~\ref{ch:problem}, The Problem and its Challenges,} states the problem and the challenges that shaped the system before it was built.
    \item \textbf{Chapter~\ref{ch:proposition}, \sysname{} Proposal,} presents the system in four parts --- indexing, retrieval, generation and evaluation --- with the decisions behind each component and the evaluation results.
    \item \textbf{Chapter~\ref{ch:prototype-demonstration}, \sysname{} in Practice,} shows the system answering real questions end to end.
    \item \textbf{Chapter~\ref{ch:conclusion}, Conclusions and Future Work,} answers the research questions, states the limitations and outlines future work.
\end{itemize}
